%% File: wzusetup.tex
%% Derived from: thusetup.tex (Tsinghua University Thesis Template)
%%
%% Major Changes:
%%   - Updated the basic information and configuration of the paper for Wenzhou University

% !TeX root = ./wzuthesis-example.tex

% 论文基本信息配置

\wzusetup{
  %******************************
  % 注意：
  %   1. 配置里面不要出现空行
  %   2. 不需要的配置信息可以删除
  %   3. 建议先阅读文档中所有关于选项的说明
  %******************************
  %
  % 输出格式
  %   选择打印版（print）或用于提交的电子版（electronic），前者会插入空白页以便直接双面打印
  %
  output = electronic,
  % 格式类型
  %   默认为论文（thesis），也可以设置为开题报告（proposal）
  % thesis-type = proposal,
  %
  % 标题
  %   可使用“\\”命令手动控制换行
  %
  clc                = {TP391.4},
  udc                = {004.8},
  student-id = {24451352011},
  secret-level = {公开},
  title  = {面向高效的无监督分割方法},
  %
  % 培养单位
  %   填写所属院系的全名
  %
  department = {计算机与人工智能学院},
  %
  % 学科
  %   1. 研究生学术型学位，获得一级学科授权的学科填写一级学科名称，其他填写二级学科名称
  %   2. 本科生填写专业名称，第二学位论文需标注“（第二学位）”
  %
  major  = {计算机科学与技术},
  %
  % 专业领域
  %   1. 设置专业领域的专业学位类别，填写相应专业领域名称
  %   2. 2019 级及之前工程硕士学位论文，在 `engineering-field` 填写相应工程领域名称
  %   3. 其他专业学位类别的学位论文无需此信息
  %
  % professional-field  = {计算机技术},
  % professional-field* = {Computer Technology},
  %
  % 姓名
  %
   author  = {冯兴宇},
  %
  % 学号
  % 仅当书写开题报告时需要（同时设置 `thesis-type = proposal'）
  %
  % student-id = {2000310000},
  %
  % 指导教师
  %   中文姓名和职称之间以英文逗号“,”分开，下同
  %
  supervisor  = {李洪, 教授},
  %
  % 副指导教师
  %
%  associate-supervisor  = {陈文光, 教授},
  %
  % 联合指导教师
  %
  % co-supervisor  = {某某某, 教授},
  % co-supervisor* = {Professor Mou Moumou},
  %
  % 日期
  %   使用 ISO 格式；默认为当前时间
  %
  % date = {2019-07-07},
  %
  % 是否在中文封面后的空白页生成书脊（默认 false）
  %
  finish-date = {2026年6月16日},
  include-spine = false,
  %
  % 密级和年限
  %   秘密, 机密, 绝密
  %
  % secret-level = {秘密},
  % secret-year  = {10},
  %
  % 博士后专有部分
  %
  % clc                = {分类号},
  % udc                = {UDC},
  % id                 = {编号},
  % discipline-level-1 = {计算机科学与技术},  % 流动站（一级学科）名称
  % discipline-level-2 = {系统结构},          % 专业（二级学科）名称
  % start-date         = {2011-07-01},        % 研究工作起始时间
}

% 载入所需的宏包

% 定理类环境宏包
\usepackage{amsthm}
% 也可以使用 ntheorem
% \usepackage[amsmath,thmmarks,hyperref]{ntheorem}

% 可以使用 nomencl 生成符号和缩略语说明
% \usepackage{nomencl}
% \makenomenclature

% 表格加脚注
\usepackage{threeparttable}

% 表格中支持跨行
\usepackage{multirow}

% 固定宽度的表格。
% \usepackage{tabularx}

% 跨页表格
\usepackage{longtable}

% 算法
\usepackage{algorithm}
\usepackage{algorithmic}

% 量和单位
\usepackage{siunitx}

% 参考文献使用 BibTeX + natbib 宏包
% 顺序编码制
\usepackage[sort]{natbib}
\bibliographystyle{wzuthesis-numeric}

% 著者-出版年制
% \usepackage{natbib}
% \bibliographystyle{wzuthesis-author-year}

% 生命科学学院要求使用 Cell 参考文献格式（2023 年以前使用 author-date 格式）
% \usepackage{natbib}
% \bibliographystyle{cell}

% 本科生参考文献的著录格式
% \usepackage[sort]{natbib}
% \bibliographystyle{wzuthesis-bachelor}

% 参考文献使用 BibLaTeX 宏包
% \usepackage[style=wzuthesis-numeric]{biblatex}
% \usepackage[style=wzuthesis-author-year]{biblatex}
% \usepackage[style=gb7714-2015]{biblatex}
% \usepackage[style=apa]{biblatex}
% \usepackage[style=mla-new]{biblatex}
% 声明 BibLaTeX 的数据库
% \addbibresource{ref/refs.bib}

% 定义所有的图片文件在 figures 子目录下
\graphicspath{{figures/}}

% 数学命令
\makeatletter
\newcommand\dif{%  % 微分符号
  \mathop{}\!%
  \ifwzu@math@style@TeX
    d%
  \else
    \mathrm{d}%
  \fi
}
\makeatother

% hyperref 宏包在最后调用
\usepackage{hyperref}

% 中英文图/表标题
\usepackage[list=off]{bicaption} % 禁用 bicaption 自动添加到 lof
\captionsetup[figure][bi-second]{name=Figure}
\captionsetup[table][bi-second]{name=Table}